%% =====================================================================
%% Accessibility: what screen readers get in the tagged PDF, beyond LaTeX's own tagging. Nothing here
%% changes the look of the pages. Loaded at the end of packages.tex.
%%   1. Words in maths: operator names (sin, tr, arg max) and \mathword{Re} are read as one word
%%   2. Tables: the first row of a table is its header (so each cell is read with its column heading)
%%   3. Formulas with an alternative text of your own: \formulaalt{...}
%%   4. Acronyms: the short form carries its expansion ("BEM" is read as "blade element momentum")
%%   5. Units read as words: \qty{10}{\metre\per\second} is read as "10 metres per second"
%%   6. Other languages: text in an otherlanguage environment is marked with its language (Dutch)
%% =====================================================================
\makeatletter
\ExplSyntaxOn

%% ---------------------------------------------------------------------
%% 1. Words in maths. In the MathML every letter of "sin" or "Re" is a variable of its own, which screen
%% readers spell out ("s i n"). Operator names (\sin, \max, \DeclareMathOperator, \operatorname) and
%% \mathword{...} get the whole word as one identifier. \mathword{Re} looks like \mathrm{Re}: use it for
%% names of more than one letter, such as dimensionless numbers (commands.tex).
%% ---------------------------------------------------------------------
\text_declare_purify_equivalent:Nn \vphantom { \use_none:n }
\text_declare_purify_equivalent:Nn \, { ~ }
\text_declare_purify_equivalent:Nn \; { ~ }
\cs_new:Npn \__thesis_mathword:n #1 { \lua_escape:e { \text_purify:n {#1} } }
\NewDocumentCommand \mathword { m } { \mathrm {#1} }
% inside: no structure elements of its own (a \vphantom in \max would leave one without a parent, as the
% word replaces what is inside)
\cs_new_eq:NN \__thesis_mathword_nostruct: \thesis_luamml_no_structure:
% (luamml, which makes the MathML, is loaded at the start of the document)
\AddToHook { begindocument }
  {
    \cs_if_exist:NT \luamml_annotate:nen
      {
        % amsmath's \qopname, which all operator names go through, with the name as one identifier
        \cs_set_protected:Npn \qopname #1#2#3
          {
            \luamml_annotate:nen { 1 }
              { nucleus = true, mathml = { [0] = 'mi', '\__thesis_mathword:n {#3}' } }
              { \mathop { \__thesis_mathword_nostruct: #1 \kern \c_zero_dim \use:c { operator@font } #3 } \use:c { n#2limits@ } }
          }
        \RenewDocumentCommand \mathword { m }
          {
            \luamml_annotate:nen { 1 }
              { nucleus = true, mathml = { [0] = 'mi', '\__thesis_mathword:n {#1}' } }
              { \mathord { \__thesis_mathword_nostruct: \mathrm {#1} } }
          }
      }
  }

%% \layouttable before a tabular that only places text (the committee, the colophon, the curriculum
%% vitae): its rows and cells are then tagged as text blocks, not as a table, so a screen reader reads
%% them as text and checkers do not ask for header cells. It holds until the end of the group.
\NewDocumentCommand { \layouttable } { }
  { \tag_if_active:T { \tagpdfsetup { table/tagging = div } } }

\tag_if_active:T
  {
%% ---------------------------------------------------------------------
%% 2. Tables. The first row of every table in a table float (and of xltabular and longtable) is marked as
%% its header row. A table with a header of two rows: \tagpdfsetup{table/header-rows={1,2}} in the
%% table, before the tabular; none: \tagpdfsetup{table/header-rows={}}. Tables used only to place text:
%% \layouttable (above).
%% ---------------------------------------------------------------------
    \AddToHook { env/table/begin } { \tagpdfsetup { table/header-rows = 1 } }
    \AddToHook { env/sidewaystable/begin } { \tagpdfsetup { table/header-rows = 1 } }
    \AddToHook { env/xltabular/begin } { \tagpdfsetup { table/header-rows = 1 } }
    \AddToHook { env/longtable/begin } { \tagpdfsetup { table/header-rows = 1 } }

%% A \makecell (or \thead) is a small tabular inside a cell: set as a layout table, so that a header that
%% breaks over two lines stays one header cell with text, not a table inside a header cell.
    \AddToHook { begindocument/before }
      {
        \cs_if_exist:NT \mcell@@@tabular
          {
            \cs_new_eq:NN \__thesis_mcell_tabular:w \mcell@@@tabular
            \cs_set:Npn \mcell@@@tabular #1#2 \@nil #3
              { \group_begin: \layouttable \__thesis_mcell_tabular:w {#1} #2 \@nil {#3} \group_end: }
          }
      }

%% Links: each link annotation also gets a short description (/Contents: "Go to destination ...", or the
%% web address). PDF/UA-2 does not ask for it, but checkers built on PDF/UA-1 (PAC, Acrobat) do, and
%% LaTeX adds it by itself only for PDF/UA-1.
    \AddToHook { begindocument }
      {
        \clist_map_inline:nn { GoTo , GoToR , URI }
          {
            \socket_if_exist:nT { hyp/link/#1/Contents }
              { \socket_assign_plug:nn { hyp/link/#1/Contents } { default } }
          }
      }

%% ---------------------------------------------------------------------
%% 3. Formulas: an alternative text of your own, for a formula whose MathML reads badly (a large matrix,
%% a diagram drawn in maths). Write it right before the formula:
%%   \formulaalt{The stiffness matrix: a block matrix of four blocks, ...}
%%   \begin{equation} ... \end{equation}
%% Screen readers then read that text instead of the formula.
%% ---------------------------------------------------------------------
    \tl_new:N \g__thesis_formula_alt_tl
    \NewDocumentCommand \formulaalt { m } { \tl_gset:Nn \g__thesis_formula_alt_tl {#1} }
    \NewTaggingSocketPlug { math/content } { thesis }
      {
        % as LaTeX's alt+source, which the PDF/UA-2 set-up uses
        \tl_set:Ne \l__math_content_alt_tl { \l__math_content_template_tl }
        \tl_set:Ne \l__math_content_AF_source_tl { \l__math_texsource_template_tl }
        \tl_set:Nn \l__math_content_AF_mathml_tl { }
        \tl_set:Nn \l__math_content_actual_tl { }
        \tl_if_empty:NF \g__thesis_formula_alt_tl
          {
            \tl_set:Ne \l__math_content_alt_tl { \text_purify:n { \g__thesis_formula_alt_tl } }
            \bool_set_true:N \l__tag_math_alt_bool
            \tl_gclear:N \g__thesis_formula_alt_tl
          }
      }
    \AssignTaggingSocketPlug { math/content } { thesis }

%% ---------------------------------------------------------------------
%% 4. Acronyms: \gls{bem} after its first use gives "BEM", marked with its expansion (/E), so a screen
%% reader can say "blade element momentum theory". (formatNomenclature.tex calls \thesisexpansion.)
%% ---------------------------------------------------------------------
    \cs_gset_protected:Npn \thesisexpansion #1#2 % expansion, text
      {
        \mode_if_horizontal:F { \leavevmode }
        \tag_mc_end_push:
        \use:e { \exp_not:N \tag_struct_begin:n { tag = Span , E = { \text_purify:n {#1} } } }
        \tag_mc_begin:n { }
        #2
        \tag_mc_end:
        \tag_struct_end:
        \tag_mc_begin_pop:n { }
      }

%% ---------------------------------------------------------------------
%% 5. Units read as words: \unit{\metre\per\second} carries the expansion "metres per second", and
%% \qty{10}{\metre\per\second} "10 metres per second" (singular after 1). Only units written with
%% siunitx's macros (\kilo\metre, \per, \squared, ...) are spelt out; a unit written as text, such as
%% \unit{m/s}, is left as it is. Units of your own: \thesisunitword{\furlong}{furlong}.
%% ---------------------------------------------------------------------
    \prop_new:N \g__thesis_unit_words_prop
    \prop_new:N \g__thesis_unit_prefix_prop
    \cs_gset_protected:Npn \thesisunitword #1#2
      { \prop_gput:Nen \g__thesis_unit_words_prop { \cs_to_str:N #1 } {#2} }
    \cs_generate_variant:Nn \prop_gput:Nnn { Nen }
    \cs_gset_protected:Npn \__thesis_unit_pair:w #1 / #2 \q_stop #3 { \prop_gput:Nnn #3 {#1} {#2} }
    \clist_map_inline:nn
      {
        metre/metre, meter/metre, second/second, gram/gram, kilogram/kilogram, kelvin/kelvin,
        ampere/ampere, mole/mole, candela/candela, newton/newton, pascal/pascal, joule/joule,
        watt/watt, hertz/hertz, volt/volt, ohm/ohm, coulomb/coulomb, farad/farad, henry/henry,
        tesla/tesla, weber/weber, siemens/siemens, radian/radian, steradian/steradian,
        degreeCelsius/degree~Celsius, celsius/degree~Celsius, degree/degree, arcminute/arcminute,
        arcsecond/arcsecond, litre/litre, liter/litre, hour/hour, minute/minute, day/day,
        tonne/tonne, bar/bar, percent/percent, electronvolt/electronvolt, decibel/decibel,
        neper/neper, bel/bel, hectare/hectare, astronomicalunit/astronomical~unit,
        dalton/dalton, knot/knot, nauticalmile/nautical~mile, angstrom/angstrom, lux/lux,
        lumen/lumen, becquerel/becquerel, gray/gray, sievert/sievert, katal/katal
      }
      { \__thesis_unit_pair:w #1 \q_stop \g__thesis_unit_words_prop }
    \clist_map_inline:nn
      {
        quecto/quecto, ronto/ronto, yocto/yocto, zepto/zepto, atto/atto, femto/femto, pico/pico,
        nano/nano, micro/micro, milli/milli, centi/centi, deci/deci, deca/deca, deka/deca,
        hecto/hecto, kilo/kilo, mega/mega, giga/giga, tera/tera, peta/peta, exa/exa,
        zetta/zetta, yotta/yotta, ronna/ronna, quetta/quetta
      }
      { \__thesis_unit_pair:w #1 \q_stop \g__thesis_unit_prefix_prop }
    \seq_new:N \l__thesis_unit_seq
    \tl_new:N \l__thesis_unit_tmp_tl
    \tl_new:N \l__thesis_unit_words_tl
    \tl_new:N \l__thesis_unit_number_tl
    \bool_new:N \l__thesis_unit_ok_bool
    \cs_generate_variant:Nn \prop_get:NnNTF { NV }
    \cs_generate_variant:Nn \regex_match:nnTF { nV }
    \cs_generate_variant:Nn \str_case:nnF { V }
    % each unit is kept as {square or cubic}{prefix and unit}{squared or cubed}; "per" as {}{}{}
    \cs_gset_protected:Npn \__thesis_unit_words:nn #1#2 % unit, plural
      {
        \seq_clear:N \l__thesis_unit_seq
        \bool_set_true:N \l__thesis_unit_ok_bool
        \tl_clear:N \l_tmpa_tl % a prefix waiting for its unit
        \tl_clear:N \l_tmpb_tl % square or cubic, waiting
        \tl_map_inline:nn {#1}
          {
            \token_if_cs:NTF ##1
              {
                \str_set:Ne \l_tmpa_str { \cs_to_str:N ##1 }
                \prop_get:NVNTF \g__thesis_unit_prefix_prop \l_tmpa_str \l__thesis_unit_tmp_tl
                  { \tl_set_eq:NN \l_tmpa_tl \l__thesis_unit_tmp_tl }
                  {
                    \prop_get:NVNTF \g__thesis_unit_words_prop \l_tmpa_str \l__thesis_unit_tmp_tl
                      {
                        \seq_put_right:Ne \l__thesis_unit_seq
                          { { \exp_not:V \l_tmpb_tl } { \exp_not:V \l_tmpa_tl \exp_not:V \l__thesis_unit_tmp_tl } { } }
                        \tl_clear:N \l_tmpa_tl \tl_clear:N \l_tmpb_tl
                      }
                      {
                        \str_case:VnF \l_tmpa_str
                          {
                            { per } { \seq_put_right:Nn \l__thesis_unit_seq { { } { } { } } }
                            { square } { \tl_set:Nn \l_tmpb_tl { square } }
                            { cubic } { \tl_set:Nn \l_tmpb_tl { cubic } }
                            { squared } { \__thesis_unit_suffix:n { squared } }
                            { cubed } { \__thesis_unit_suffix:n { cubed } }
                          }
                          { \bool_set_false:N \l__thesis_unit_ok_bool }
                      }
                  }
              }
              { \tl_if_blank:nF {##1} { \bool_set_false:N \l__thesis_unit_ok_bool } }
          }
        \tl_clear:N \l__thesis_unit_words_tl
        \bool_lazy_and:nnT { \l__thesis_unit_ok_bool } { ! \seq_if_empty_p:N \l__thesis_unit_seq }
          {
            % plural: the last unit before the first "per"
            \int_zero:N \l_tmpa_int \int_zero:N \l_tmpb_int
            \bool_if:nT {#2}
              {
                \seq_map_indexed_inline:Nn \l__thesis_unit_seq
                  {
                    \tl_if_blank:eTF { \use_ii:nnn ##2 }
                      { \int_compare:nNnT { \l_tmpb_int } = { 0 } { \int_set:Nn \l_tmpb_int {##1} } }
                      { \int_compare:nNnT { \l_tmpb_int } = { 0 } { \int_set:Nn \l_tmpa_int {##1} } }
                  }
              }
            \seq_map_indexed_inline:Nn \l__thesis_unit_seq
              { \__thesis_unit_join:nnnn {##1} ##2 }
          }
      }
    \cs_gset_protected:Npn \__thesis_unit_join:nnnn #1#2#3#4 % index, power, unit, suffix
      {
        \tl_put_right:Ne \l__thesis_unit_words_tl
          {
            \tl_if_empty:NF \l__thesis_unit_words_tl { ~ }
            \tl_if_blank:nTF {#3}
              { per }
              {
                \tl_if_blank:nF {#2} { #2 ~ }
                \int_compare:nNnTF {#1} = { \l_tmpa_int } { \__thesis_unit_s:n {#3} } {#3}
                \tl_if_blank:nF {#4} { ~ #4 }
              }
          }
      }
    \cs_gset_protected:Npn \__thesis_unit_suffix:n #1
      {
        \seq_pop_right:NNTF \l__thesis_unit_seq \l__thesis_unit_tmp_tl
          { \seq_put_right:Ne \l__thesis_unit_seq { \exp_last_unbraced:NV \__thesis_unit_setsuffix:nnnn \l__thesis_unit_tmp_tl {#1} } }
          { \bool_set_false:N \l__thesis_unit_ok_bool }
      }
    \cs_gset:Npn \__thesis_unit_setsuffix:nnnn #1#2#3#4 { { \exp_not:n {#1} } { \exp_not:n {#2} } { \exp_not:n {#4} } }
    % the plural of a unit word: an s, except for a few ("degree Celsius" -> "degrees Celsius")
    \cs_gset:Npn \__thesis_unit_s:n #1
      {
        \str_case:nnF {#1}
          {
            { hertz } {#1} { siemens } {#1} { lux } {#1} { percent } {#1}
            { degree~Celsius } { degrees~Celsius }
          }
          { #1 s }
      }
    % the wrappers: in text only (in maths the MathML has the units)
    \AddToHook { begindocument }
      {
        \NewCommandCopy \__thesis_qty:w \qty % (copies that keep working once \qty is redefined)
        \NewCommandCopy \__thesis_unit:w \unit
        \RenewDocumentCommand \qty { O{} m m }
          {
            \mode_if_math:TF { \__thesis_qty:w [#1] {#2} {#3} }
              {
                \tl_set:Ne \l__thesis_unit_number_tl { \tl_trim_spaces:n {#2} }
                \bool_set:Nn \l_tmpa_bool { ! \str_if_eq_p:Vn \l__thesis_unit_number_tl { 1 } }
                \__thesis_unit_words:nn {#3} { \l_tmpa_bool }
                % (only a plain number: one in another form is read as written)
                \regex_match:nVTF { \A[+\-]?[0-9]+([.,][0-9]+)?\Z } \l__thesis_unit_number_tl
                  {
                    \tl_if_empty:NTF \l__thesis_unit_words_tl
                      { \__thesis_qty:w [#1] {#2} {#3} }
                      { \thesisexpansion { #2 ~ \l__thesis_unit_words_tl } { \__thesis_qty:w [#1] {#2} {#3} } }
                  }
                  { \__thesis_qty:w [#1] {#2} {#3} }
              }
          }
        \RenewDocumentCommand \unit { O{} m }
          {
            \mode_if_math:TF { \__thesis_unit:w [#1] {#2} }
              {
                \__thesis_unit_words:nn {#2} { \c_true_bool }
                \tl_if_empty:NTF \l__thesis_unit_words_tl
                  { \__thesis_unit:w [#1] {#2} }
                  { \thesisexpansion { \l__thesis_unit_words_tl } { \__thesis_unit:w [#1] {#2} } }
              }
          }
      }

%% ---------------------------------------------------------------------
%% 6. Other languages: an otherlanguage environment (babel) is a division marked with its language, e.g.
%% nl for Dutch, and \foreignlanguage a span, so screen readers switch pronunciation. (A switch such as
%% \selectlanguage, or otherlanguage*, has no end to close a structure at: use otherlanguage instead.)
%% ---------------------------------------------------------------------
    \AddToHook { begindocument/before }
      {
        \cs_set_eq:NN \__thesis_otherlanguage:n \otherlanguage
        \cs_set_eq:NN \__thesis_endotherlanguage: \endotherlanguage
        \cs_set_protected:Npn \otherlanguage #1
          {
            \__thesis_otherlanguage:n {#1}
            \par
            \tag_struct_begin:n { tag = Div , lang = \localeinfo* { tag.bcp47 } }
          }
        \cs_set_protected:Npn \endotherlanguage
          { \par \tag_struct_end: \__thesis_endotherlanguage: }
        % a few words in another language, \foreignlanguage{dutch}{proefschrift}: a span with that language
        \NewCommandCopy \__thesis_foreignlanguage:w \foreignlanguage
        \RenewDocumentCommand \foreignlanguage { O{} m m }
          { \__thesis_foreignlanguage:w [#1] {#2} { \__thesis_langspan:n {#3} } }
      }
    \cs_new_protected:Npn \__thesis_langspan:n #1
      {
        \mode_if_horizontal:F { \leavevmode }
        \tag_mc_end_push:
        \use:e { \exp_not:N \tag_struct_begin:n { tag = Span , lang = \localeinfo* { tag.bcp47 } } }
        \tag_mc_begin:n { }
        #1
        \tag_mc_end:
        \tag_struct_end:
        \tag_mc_begin_pop:n { }
      }
  }
\cs_if_exist:NF \thesisexpansion { \cs_new_protected:Npn \thesisexpansion #1#2 {#2} }

\ExplSyntaxOff
\makeatother
